% =====================================================================
% SECCION 3 -- PLANTEAMIENTO DEL PROBLEMA Y ANALISIS DEL CONTEXTO
%
% Creada el 14/09/2026 al reestructurar el informe al indice del modelo
% del curso (NeuroSegment). Orden del modelo:
%   3.1 Descripcion del Caso
%   3.2 Identificacion de Variables y Restricciones
%   3.3 Restricciones del Proyecto (3.3.1 Tecnicas, 3.3.2 Eticas, 3.3.3 Legales)
%   3.4 Analisis Tecnico, Etico, Social, Ambiental
%
% ORIGEN DEL TEXTO
% - 3.1, 3.2 y el primer parrafo de 3.3.1 son texto de la antigua
%   Introduccion (otro integrante), trasladado sin cambiar sus palabras
%   salvo la remision "Seccion de Metodologia" -> \ref{sec:diseno}.
% - La frase que abre 3.2 y el segundo parrafo de 3.3.1 son de integracion
%   (frente etico-legal), redactados con datos que ya constaban en el informe.
% - 3.3.2, 3.3.3 y 3.4 NO estan aqui: se generan desde
%   secciones/03_marco_normativo.md y secciones/03_analisis_impacto.md,
%   y main.tex los incluye justo despues de este archivo.
% =====================================================================
\section{Planteamiento del Problema y Análisis del Contexto}
\label{sec:planteamiento}

\subsection{Descripción del Caso}
\label{sec:contexto}

El Aeropuerto Internacional Jorge Chávez, operado por Lima Airport
Partners (LAP), no distribuye a sus pasajeros de manera uniforme. En
determinadas franjas horarias, y en puntos concretos del terminal
(filtros de seguridad, control migratorio, mostradores de
\emph{check-in}, salas de embarque), se forman concentraciones muy
superiores a las del resto del recinto en ese mismo momento.

La aglomeración, por sí sola, no constituye el problema. Lo constituye
cuando la cantidad de personas reunidas en un punto excede, durante un
tiempo sostenido, la capacidad de atención dispuesta para ese punto: es
entonces cuando aparecen las colas que se alargan, los tiempos de espera
que comprometen la conexión de un vuelo y las condiciones de
incomodidad o de riesgo en un espacio cerrado. La magnitud que interesa
medir no es, por tanto, el volumen total de pasajeros del terminal, sino
una combinación de tres variables: \emph{cuántas} personas, en
\emph{qué} zona y durante \emph{cuánto} tiempo.

Medir esa combinación importa porque las decisiones que la atienden son
de corto plazo y reversibles: habilitar un mostrador adicional,
redirigir el flujo hacia un filtro descargado, adelantar un relevo de
personal o abrir una segunda fila. Todas ellas dependen de conocer dónde
se está formando la concentración \emph{mientras} se forma. Un recuento
que llega al cierre de la jornada sirve para planificar el mes
siguiente, pero no para resolver la cola de esta mañana.

El procedimiento con que LAP resuelve hoy esa decisión no consta en la
documentación disponible para el equipo, de modo que este informe no lo
describe ni lo evalúa.\verificar{cómo se decide hoy la dotación de
personal y la apertura de mostradores en el terminal: procedimiento,
periodicidad y responsable. Solicitado a LAP como punto C-08 de
PENDIENTES.md.} Sí puede enunciarse la limitación genérica de cualquier
método basado en observación humana periódica, que es la alternativa
frente a la cual se plantea este proyecto: cubre un punto a la vez,
produce una lectura discontinua en el tiempo, y su registro depende de
que alguien lo anote, de modo que no deja una serie homogénea sobre la
que puedan compararse dos días, dos zonas o dos franjas horarias.

Conviene precisar el alcance de todo lo anterior para no sobredimensionar
el problema que el proyecto atiende. El equipo no cuenta, al momento de
escribir este informe, con evidencia documentada por LAP que cuantifique
el problema en el propio terminal: no dispone de incidentes de
aglomeración registrados, de tiempos de cola medidos ni de quejas de
pasajeros contabilizadas.\verificar{incidentes por aglomeración
documentados, tiempos de cola observados y quejas recibidas. Es el
pedido más importante del análisis: sin él, el examen de
proporcionalidad de la Sección 3 se sostiene sobre una premisa no
acreditada. Solicitado a LAP como punto C-07 de PENDIENTES.md.}
Tampoco consta el aforo del terminal ni su superficie por
zona.\verificar{aforo declarado y superficie de las zonas del terminal.
No figura ninguna de las dos magnitudes en la documentación del
proyecto.} El proyecto parte, en consecuencia, de una necesidad
operativa plausible y compartida por la literatura de conteo de
multitudes en espacios de alto tránsito, y no de una cifra propia de
incidentes en el Aeropuerto Jorge Chávez. Todas las magnitudes que este
informe sí reporta proceden de conjuntos de datos públicos o de
grabaciones propias del equipo, y se identifican como tales en cada
caso.

\subsection{Identificación de Variables y Restricciones}
\label{sec:alcance}

Las variables que el sistema debe estimar son las tres enunciadas en la
Sección~\ref{sec:contexto}: cuántas personas hay, en qué zona y durante
cuánto tiempo. Su delimitación depende de las condiciones siguientes.

\textbf{Alcance funcional.} El proyecto cubre la estimación de la
concentración de personas por zona a partir de imágenes de cámaras fijas.
Queda expresamente fuera de su alcance cualquier forma de identificación,
reidentificación o caracterización individual de las personas detectadas:
el sistema está pensado para responder \emph{cuántas personas hay, y
dónde}, no \emph{quién es cada una}. La Sección~\ref{sec:ant-p2pnet}
señala, sin embargo, que esta distinción es una cuestión de política de
persistencia y no solo de algoritmo: el mismo modelo que localiza personas
puede, si sus salidas se retienen y se asocian entre cuadros, aproximarse
a un seguimiento de trayectorias. Por eso la granularidad de salida que
finalmente se implemente es, como se indica en el objetivo específico 3,
una decisión que debe quedar documentada por escrito antes de cerrar la
Sección~\ref{sec:diseno}.

\textbf{Zonas cubiertas.} Las zonas concretas del terminal sobre las que
se entrena y evalúa el sistema dependen de qué cámaras y qué material
audiovisual habilite LAP para el proyecto, dato que a la fecha de este
informe no ha sido confirmado y que solo LAP puede proporcionar. Donde este informe menciona zonas a título de ejemplo
(filtros de seguridad, mostradores de \emph{check-in}, salas de embarque),
lo hace como ilustración del tipo de espacio al que el sistema podría
aplicarse, no como una cobertura ya acordada con el cliente. El número de
zonas que finalmente se defina, y el número de cámaras que las
cubra, tampoco están fijados.\verificar{número de zonas del terminal a
cubrir y número de cámaras disponibles por zona. El prototipo trabaja
hoy con dos cámaras sobre una única zona, pero es material propio del
equipo grabado en la universidad, no del terminal: no es una cifra
extrapolable.}

\subsection{Restricciones del Proyecto}
\label{sec:restricciones}

\subsubsection{Restricciones Técnicas}

\textbf{Fase del proyecto y limitación de dominio.} Esta primera entrega
corresponde a Capstone I: un desarrollo y validación conceptual del
sistema sobre conjuntos de datos públicos de conteo de multitudes, sin
acceso todavía a material propio del Aeropuerto Jorge Chávez. Esto impone
la limitación más importante del proyecto en esta etapa, y conviene
enunciarla sin atenuarla porque condiciona la lectura de todo lo que
sigue: como documenta en detalle la Sección~\ref{sec:ant-p2r}, los propios
autores de la literatura revisada reconocen una brecha de dominio entre
los conjuntos públicos de entrenamiento y un entorno de despliegue nuevo,
y señalan que el desempeño se degrada de forma marcada cuando no media una
etapa de adaptación de dominio. En consecuencia, ninguna cifra de
desempeño que se reporte en la Sección de Resultados debe leerse como una
predicción de exactitud para el Aeropuerto Jorge Chávez: es, en el mejor
de los casos, una cota de referencia sobre datos públicos, sujeta a
validarse en sitio antes de comunicarse a LAP como una expectativa de
desempeño (Sección~\ref{sec:auditoria}).

A esa limitación se suman dos restricciones propias del entorno. El
sistema debe operar sobre las cámaras fijas ya instaladas, que no
comparten resolución, altura ni ángulo. Y buena parte de la jornada el
terminal opera en baja densidad, régimen que, según reconocen los propios
autores de CSRNet, no es el terreno donde su método destaca
(Sección~\ref{sec:ant-csrnet}).
